In incident response, command belongs to a person: the incident commander, appointed by
the organization. The commander works from a vehicle or a handheld that is often cut off,
and the organization may then move command to a deputy or a superior, while the devices of
the people involved cannot reach one another and crews keep acting on the orders those
devices record. We ask how an intermittently connected incident information system can
represent such a succession, enforce it on devices that cannot reach each other, and
reconcile it when they reconnect, and what each succession policy costs the people who act
on its orders. We separate a person's authority, conferred by signed grants and by
succession rules the organization fixes in advance, from a device's enforcement of it, and
give every order one of four statuses that crews and commanders can see. On an executable
model of a forest-fire incident, we compare six succession policies across three partition
cases and three states of the commander, over \tbTotalRuns{} simulated incidents. A
commander's device that stops on its own after losing contact with the deputy can be
revived by a stale heartbeat relayed through another vehicle; a timeout acknowledged by
the deputy prevents this with a margin of only the clock error. A takeover that takes
effect when it is signed, rather than when the successor's device accepts it, turns orders
given in good faith into residual orders after the fact. Settlement of orders is safe, but
it completes only once the predecessor's records arrive or the appointing officer acts.
\pending{Prototype runs on vehicle edges and an expert walkthrough with rescue-service
officers: planned.}
